\reading{LTR-5}{Liquidity and Reserves Management: Strategies and Policies}{DEFER}{0}
\markboth{LTR-5\ \ Liquidity and Reserves Management}{}

\begin{lrkeybox}[Learning objectives — official 2026 spine wording]
\begin{enumerate}[label=\textbf{\alph*.}]
\item Calculate a bank's net liquidity position and explain factors that affect the supply and demand of liquidity at a bank.
\item Compare strategies that a bank can use to meet demands for additional liquidity.
\item Estimate a bank's liquidity needs through three methods (sources and uses of funds, structure of funds, and liquidity indicators).
\end{enumerate}
{\footnotesize\color{FRMGrey}Source: Peter Rose and Sylvia Hudgins, \textit{Bank Management \& Financial Services}, 9th Edition, Chapter 11.}
\end{lrkeybox}

\begin{lrtrapbox}[Two objectives in the source notes are off-syllabus]
The Crash layer carries two further objectives: \texttt{68.d}, summarizing how a US bank calculates its legal reserves, and \texttt{68.e}, differentiating factors affecting the choice among alternate sources of reserves. Neither appears in the 2026 spine, which gives this reading exactly three objectives. Both have been \textbf{excluded} rather than carried with a caveat. If a question bank or tier document drills legal reserve calculation or reserve-source selection under this reading, that material is dead weight.
\end{lrtrapbox}

% =====================================================================
\lo{LTR-5 a}{Calculate a bank's net liquidity position and explain factors that affect the supply and demand of liquidity at a bank.}

\begin{lrfmlbox}[Net liquidity position]
\[
L_t \;=\;
\underbrace{\!\begin{aligned}
&\text{Incoming deposits} + \text{Revenues from the sale of nondeposit services}\\
&+ \text{Customer loan repayments} + \text{Sales of assets} + \text{Borrowings from the money market}
\end{aligned}}_{\textstyle\text{Supplies of liquidity flowing \emph{into} the financial firm}}
\]
\vspace{-6pt}
\[
\;-\;
\underbrace{\!\begin{aligned}
&\text{Deposit withdrawals} + \text{Volume of acceptable loan requests} + \text{Repayments of borrowings}\\
&+ \text{Other operating expenses} + \text{Dividend payments to stockholders}
\end{aligned}}_{\textstyle\text{Demands \emph{on} the financial firm for liquidity}}
\]
\vk{$L_t$ = the financial firm's net liquidity position at time $t$. Five supply terms, five demand terms. A positive $L_t$ is a liquidity surplus; a negative $L_t$ is a liquidity deficit.}
\end{lrfmlbox}

\subsection{Demands for liquidity}

\begin{enumerate}
\item Customers withdrawing money from their accounts.
\item Credit requests from customers the institution wishes to keep, either as new loan requests or drawings upon existing credit lines.
\item Paying off previous borrowings, such as loans the institution may have received from other financial firms or from the central bank.
\item Operating expenses, such as salary.
\item Payment of income taxes or cash dividends to stockholders, which periodically gives rise to a demand for immediately spendable cash.
\end{enumerate}

\subsection{Supplies of liquidity}

\begin{enumerate}
\item Receipt of new customer deposits.
\item Customers repaying their loans.
\item Sales of assets, especially marketable securities from the investment portfolio.
\item Revenues (fee income) generated by selling nondeposit services.
\item Borrowings in the money market.
\end{enumerate}

\begin{lrtrapbox}[The one term that looks like it is on the wrong side]
\textbf{Acceptable loan requests are a \emph{demand} on liquidity}, not a supply — a loan the bank wants to make is cash going out. Meanwhile \textbf{customer loan repayments are a supply}. The two sit in the same formula and differ only in direction of the cash. This is the single most reliable polarity trap in the reading.
\end{lrtrapbox}

\begin{lrkeybox}[Liquidity has a critical time dimension]
\begin{itemize}
\item Some liquidity has near-term pressure: large CDs due to mature tomorrow.
\item Longer-term liquidity demands arise from seasonal, cyclical, and trend factors.
\item \textbf{Most liquidity problems arise from \emph{outside} the financial firm}, as a result of the activities of customers.
\end{itemize}
\end{lrkeybox}

\begin{lrkeybox}[The essence of liquidity management, in two statements]
\begin{enumerate}
\item \textbf{Rarely are demands for liquidity equal to the supply of liquidity} at any particular moment in time. The firm is continually in either a liquidity deficit or a liquidity surplus.
\item There is a \textbf{trade-off between liquidity and profitability}.
\end{enumerate}
\end{lrkeybox}

% =====================================================================
\lo{LTR-5 b}{Compare strategies that a bank can use to meet demands for additional liquidity.}

A financial firm is liquid only if it has access, at reasonable cost, to liquid funds in exactly the amounts required at precisely the time they are needed. Experienced liquidity managers have developed three strategies.

\begin{center}
\small
\begin{tabularx}{\textwidth}{@{}p{3.3cm} X X@{}}
\toprule
\rowcolor{FRMSoft}\textbf{Strategy} & \textbf{Mechanics and who uses it} & \textbf{Pros and cons} \\
\midrule
\textbf{1. Asset liquidity management} \newline (asset conversion) &
The oldest approach. When liquidity is needed, selected assets are converted into cash until all demands for cash are met. Maintains a stock of highly liquid assets. Used mainly by \textbf{smaller} institutions, which find it a \emph{less risky} approach than relying on borrowings. &
\textbf{Pro:} readily available cash. \textbf{Cons:} it is not costless. Selling assets means the loss of future earnings those assets would have generated; it weakens the appearance of the balance sheet by selling low-risk securities; and liquid assets generally carry the lowest rates of return of all assets. Also opportunity costs and transaction costs. \\
\addlinespace
\textbf{2. Borrowed liquidity management} \newline (liability management) &
Raising liquid funds through borrowings in the money market. Chosen by the \textbf{largest} financial institutions around the world. &
\textbf{Pros:} borrow only when funds are actually needed; leaves the volume and composition of the asset portfolio unchanged; comes with its own control lever, the interest rate offered to borrow funds. \textbf{Cons:} the \emph{most risky} approach to solving liquidity problems, but also carries the \emph{highest expected return}. Borrowing cost is always uncertain, adding greater uncertainty to earnings. \\
\addlinespace
\textbf{3. Balanced liquidity management} &
Most financial firms compromise by using both. Some expected demands are stored in assets (principally marketable securities); other anticipated needs are backstopped by advance arrangements for lines of credit; unexpected cash needs are typically met from near-term borrowings; longer-term needs are planned for and parked in short- and medium-term assets that provide cash as those needs arise. &
\textbf{Pros:} balances liquidity needs with asset management, avoids over-reliance on borrowing. \textbf{Cons:} requires time management and anticipating future needs. \\
\bottomrule
\end{tabularx}
\end{center}
\figcap{The discriminator is institution size paired with risk appetite: \textbf{small firms convert assets} because it is the less risky route; \textbf{large firms borrow} because they can, accepting the highest risk for the highest expected return. Swapping the two sizes is the obvious role-misassignment build.}

\begin{lrtrapbox}[Risk and return run together here]
Borrowed liquidity is described as \textbf{both} the most risky approach \textbf{and} the one carrying the highest expected return. A distractor that keeps one half and flips the other passes a careless read.
\end{lrtrapbox}

% =====================================================================
\lo{LTR-5 c}{Estimate a bank's liquidity needs through three methods (sources and uses of funds, structure of funds, and liquidity indicators).}

\subsection{Method 1 — Sources and uses of funds}

\begin{lrkeybox}[Two simple facts this method rests on]
\begin{itemize}
\item Liquidity \textbf{rises} as deposits increase and loans decrease.
\item Liquidity \textbf{declines} when deposits decrease and loans increase.
\end{itemize}
\end{lrkeybox}

The key steps are:

\begin{enumerate}
\item Loans and deposits must be forecast for a given planning period.
  \begin{itemize}
  \item Change in loans $= F_{\text{loan}}(\text{GDP, corporate earnings, money supply}, \ldots)$
  \item Change in deposits $= F_{\text{deposits}}(\text{personal income, inflation}, \ldots)$
  \end{itemize}
\item The estimated change in loans and deposits must be calculated: estimated change in deposits $-$ estimated change in loans.
\item The liquidity manager must estimate the net liquid funds' \term{surplus} or \term{deficit} — the \term{liquidity gap}.
  \begin{itemize}
  \item \textbf{Surplus}: more deposits than loans (positive gap).
  \item \textbf{Deficit}: more loans than deposits (negative gap).
  \end{itemize}
\end{enumerate}

\subsection{Method 2 — Structure of funds}

\textbf{First step.} Deposits and other funds sources are divided into categories based upon their \textbf{estimated probability of being withdrawn} — similar in spirit to the NSFR — and therefore lost to the financial firm.

\begin{center}
\small
\begin{tabularx}{\textwidth}{@{}p{4.6cm} X@{}}
\toprule
\rowcolor{FRMSoft}\textbf{Category} & \textbf{Definition} \\
\midrule
\textbf{``Hot money'' liabilities} \newline deposits and other borrowed funds &
Often called volatile liabilities. Very interest-sensitive, or funds that management is sure will be withdrawn during the current period. \\
\addlinespace
\textbf{Vulnerable funds} \newline customer deposits &
A substantial portion, perhaps \textbf{25 to 30 percent}, will probably be withdrawn sometime during the current time period. \\
\addlinespace
\textbf{Stable funds} &
Often called core deposits or core liabilities. Funds that management considers unlikely to be removed, except for a minor percentage of the total. \\
\bottomrule
\end{tabularx}
\end{center}

\textbf{Second step.} Calculate the total liquidity requirement. As an illustration, using the reserve rules:

\begin{lrfmlbox}[Total liquidity requirement — structure of funds]
\begin{align*}
\text{Total liquidity requirement} \;=\;& 0.95 \times (\text{Hot money funds} - \text{Required legal reserves}) \\
+\;& 0.30 \times (\text{Vulnerable deposits} - \text{Required legal reserves}) \\
+\;& 0.15 \times (\text{Stable funds} - \text{Required legal reserves}) \\
+\;& 1.00 \times (\text{Potential loans outstanding} - \text{Actual loans outstanding})
\end{align*}
\vk{The three liability reserve factors descend with stability: \textbf{0.95} on hot money, \textbf{0.30} on vulnerable, \textbf{0.15} on stable. Required legal reserves are usually \textbf{3\%} of each of the three categories. The fourth term carries a factor of \textbf{1.00} — the bank plans to fund the entire gap between the loans it could be asked for and the loans it has already made.}
\end{lrfmlbox}

\begin{lrtrapbox}[Three things that go wrong in this formula]
\begin{itemize}
\item \textbf{The reserve subtraction is inside each bracket}, applied before the factor, not to the total afterwards.
\item \textbf{The loan term takes a factor of 1.00, not a reduced one.} It is the only term at full weight, because a loan commitment the bank intends to honour is not a probabilistic outflow.
\item \textbf{The factors descend with volatility}: 0.95 hot, 0.30 vulnerable, 0.15 stable. Reversing the order is the polarity build, and any two of these three values will be offered.
\end{itemize}
\end{lrtrapbox}

\textbf{Last step.} A scenario analysis is applied to have a more holistic view of liquidity needs. Usually three scenarios are included:

\begin{itemize}
\item Best possible liquidity position (maximum deposits, minimum loans).
\item Liquidity position bearing the highest probability.
\item Worst possible liquidity position (minimum deposits, maximum loans).
\end{itemize}

\begin{lrfmlbox}[Expected liquidity requirement]
\[ \text{Expected liquidity requirement} \;=\; \sum_i \text{Prob}(i) \times G_i \]
\vk{$i$ indexes the scenarios (best, most likely, worst); $\text{Prob}(i)$ = the probability of outcome $i$; $G_i$ = the estimated liquidity surplus or deficit in outcome $i$. A probability-weighted average across the scenarios, not the worst case.}
\end{lrfmlbox}

\begin{lrtrapbox}[Best case pairs maximum deposits with minimum loans]
Deposits are a \emph{supply} of liquidity; loans are a \emph{demand}. So the best case is max deposits and \textbf{min} loans, and the worst is min deposits and \textbf{max} loans. Both halves move together, and a distractor pairing maximum deposits with maximum loans reads plausibly if the direction is not checked.
\end{lrtrapbox}

\subsection{Method 3 — Liquidity indicators}

Uses ratios to assess liquidity needs based on industry averages and experience.

\begin{center}
\small
\begin{tabularx}{\textwidth}{@{}X X@{}}
\toprule
\rowcolor{PaleGreen}\textbf{Higher ratio $=$ greater liquidity} & \cellcolor{PaleRed}\textbf{Higher ratio $=$ illiquidity} \\
\midrule
a) \textbf{Cash position indicator} \newline (cash $+$ deposits due from banks) / total assets &
a) \textbf{Capacity ratio} \newline (net loans and leases) / total assets \\
\addlinespace
b) \textbf{Liquid securities indicator} \newline (US government securities) / total assets &
b) \textbf{Loan commitments ratio} \newline (unused loan commitments) / total assets \\
\addlinespace
c) \textbf{Core deposit ratio} \newline (core deposits) / total assets &
c) \textbf{Pledged securities ratio} \newline (pledged securities) / total securities holdings \\
\addlinespace
d) \textbf{Hot money ratio} \newline (money market assets) / volatile liabilities &
d) \textbf{Deposit composition ratio} \newline (demand deposits) / time deposits \\
\addlinespace
e) \textbf{Net federal funds and repo position} &
e) \textbf{Deposit brokerage index} \newline (brokered deposits) / total deposits \\
\bottomrule
\end{tabularx}
\end{center}
\figcap{Ten ratios, five on each side, and the side is the whole examinable point. The left column counts things the bank \emph{holds}; the right column counts things the bank has \emph{committed, pledged or borrowed}. Two are easy to misplace: the \textbf{hot money ratio} sits on the liquidity side because its numerator is money-market \emph{assets}, and the \textbf{deposit composition ratio} sits on the illiquidity side because demand deposits can leave tomorrow while time deposits cannot.}

\begin{lrtrapbox}[The hot money ratio is not about hot money liabilities]
In Method 2, ``hot money'' names a \emph{liability} category. In Method 3, the \textbf{hot money ratio} puts money market \emph{assets} over volatile liabilities, so a higher value means \emph{more} liquidity. Same phrase, opposite side of the balance sheet, two objectives apart. Treat it as a Confusables Registry pair.
\end{lrtrapbox}

\subsection{A fourth approach named in the source notes}

\gapnote{The spine names three methods. The source notes add a fourth from the same chapter; it is retained here because it belongs to the reading, but it is not one of the three the objective enumerates.}

\textbf{Market signals / discipline approach.} The discipline of the financial marketplace approach is essentially a method where a sufficient liquidity position is gauged only by market signals indicating whether or not that level has been reached.

\begin{lrtrapbox}[Consolidated trap summary — LTR-5]
\begin{itemize}
\item \textbf{Polarity.} In $L_t$: \emph{acceptable loan requests} are a demand; \emph{customer loan repayments} are a supply.
\item \textbf{Role.} Small institutions use asset conversion (less risky); the largest use liability management (most risky, highest expected return).
\item \textbf{Calculation.} Structure-of-funds factors descend 0.95 / 0.30 / 0.15 with stability; the loan-gap term is at 1.00; legal reserves (usually 3\%) are subtracted \emph{inside} each bracket.
\item \textbf{Polarity.} Best case $=$ max deposits, min loans. Worst case $=$ min deposits, max loans.
\item \textbf{Sibling.} The \textbf{hot money ratio} (Method 3, an asset ratio, higher $=$ more liquid) versus \textbf{hot money liabilities} (Method 2, a funding category).
\item \textbf{Sibling.} Capacity ratio, loan commitments, pledged securities, deposit composition and deposit brokerage all sit on the \emph{illiquidity} side. The other five sit on the liquidity side.
\item \textbf{Scope.} Expected liquidity requirement is a \emph{probability-weighted} figure across three scenarios, not the worst-case figure.
\item \textbf{Scope.} Legal reserve calculation and the choice among alternate reserve sources are \textbf{not examinable in 2026}.
\end{itemize}
\end{lrtrapbox}
